
\documentclass{article}
\usepackage{colortbl}
\usepackage{makecell}
\usepackage{multirow}
\usepackage{supertabular}

\begin{document}

\newcounter{utterance}

\centering \large Interaction Transcript for game `hot\_air\_balloon', experiment `air\_balloon\_survival\_en\_negotiation\_easy', episode 3 with llama{-}3.3{-}70b{-}instruct{-}t1.0.
\vspace{24pt}

{ \footnotesize  \setcounter{utterance}{1}
\setlength{\tabcolsep}{0pt}
\begin{supertabular}{c@{$\;$}|p{.15\linewidth}@{}p{.15\linewidth}p{.15\linewidth}p{.15\linewidth}p{.15\linewidth}p{.15\linewidth}}
    \# & $\;$A & \multicolumn{4}{c}{Game Master} & $\;\:$B\\
    \hline

    \theutterance \stepcounter{utterance}  
    & & \multicolumn{4}{p{0.6\linewidth}}{
        \cellcolor[rgb]{0.9,0.9,0.9}{
            \makecell[{{p{\linewidth}}}]{
                \texttt{\tiny{[P1$\langle$GM]}}
                \texttt{Stai partecipando a un gioco di negoziazione collaborativa.} \\
\\ 
\texttt{Insieme a un altro partecipante, dovete accordarvi su un unico set di oggetti da conservare. Ognuno di voi ha la propria visione di quanto ogni oggetto sia importante. Non sai come l'altro partecipante valuta gli oggetti. Inoltre, ti viene fornito lo sforzo che ogni oggetto richiede.} \\
\texttt{Potete accordarvi su un set solo se lo sforzo totale degli oggetti selezionati non supera un limite condiviso:} \\
\\ 
\texttt{LIMITE = 3712} \\
\\ 
\texttt{Ecco i valori di sforzo dei singoli oggetti:} \\
\\ 
\texttt{Sforzo oggetto = \{"C08": 949, "C55": 388, "A43": 79, "A67": 585, "B42": 564, "A62": 230, "B19": 580, "C43": 84, "C73": 976, "C10": 274, "B36": 374, "A18": 913, "C61": 303, "C67": 578, "C42": 548\}} \\
\\ 
\texttt{Ecco la tua visione personale sull'importanza di ogni oggetto:} \\
\\ 
\texttt{Valori di importanza dell'oggetto = \{"C08": 138, "C55": 583, "A43": 868, "A67": 822, "B42": 783, "A62": 65, "B19": 262, "C43": 121, "C73": 508, "C10": 780, "B36": 461, "A18": 484, "C61": 668, "C67": 389, "C42": 808\}} \\
\\ 
\texttt{Obiettivo:} \\
\\ 
\texttt{Il tuo obiettivo è negoziare un set condiviso di oggetti che ti avvantaggi il più possibile (cioè, massimizzi l'importanza totale per TE), rimanendo entro il LIMITEE. Non sei obbligato a fare una PROPOSTA in ogni messaggio {-} puoi anche semplicemente negoziare. Tutte le tattiche sono permesse!} \\
\\ 
\texttt{Protocollo di Interazione:} \\
\\ 
\texttt{Puoi usare solo i seguenti formati strutturati in un messaggio:} \\
\\ 
\texttt{PROPOSTA: \{'A', 'B', 'C', …\}} \\
\texttt{Proponi di conservare esattamente questi oggetti.} \\
\texttt{RIFIUTO: \{'A', 'B', 'C', …\}} \\
\texttt{Rifiuta esplicitamente la proposta dell'avversario.} \\
\texttt{ARGOMENTAZIONE: \{'...'\}} \\
\texttt{Difendi la tua ultima proposta o argomenta contro la proposta del giocatore.} \\
\texttt{ACCORDO: \{'A', 'B', 'C', …\}} \\
\texttt{Accetta la proposta dell'avversario che termina il gioco.} \\
\texttt{RAGIONAMENTO STRATEGICO: \{'...'\}} \\
\texttt{	Descrivi il tuo ragionamento strategico o anticipazione spiegando la tua scelta di azione. Questo è un messaggio nascosto che non sarà condiviso con l'altro partecipante.} \\
\\ 
\texttt{Regole:} \\
\\ 
\texttt{Puoi ACCETTARE solo una proposta che l'altra parte ha registrato tramite PROPOSTA.} \\
\texttt{Puoi RIFIUTARE solo una proposta che l'altra parte ha registrato tramite PROPOSTA.} \\
\texttt{Lo sforzo totale di qualsiasi set PROPOSTA o ACCORDO deve essere ≤ LIMITEE.} \\
\texttt{NON rivelare i tuoi punteggi di importanza nascosti.} \\
\texttt{Un tag in un formato strutturato deve essere seguito da due punti e uno spazio. L'argomento deve essere un set python contenente 0 o più stringhe.} \\
\texttt{Quindi, deve essere della forma TAG: \{...\}} \\
\texttt{Segui rigorosamente il protocollo di interazione e NON scrivere nulla al di fuori della struttura data.} \\
\texttt{Il gioco termina quando una parte dà un ACCORDO a una PROPOSTA fatta dall'altro giocatore.} \\
\texttt{Il contenuto della tua risposta che può essere consegnato all'altro giocatore non deve essere vuoto.} \\
\texttt{Sono attive solo le proposte registrate tramite la struttura del formato PROPOSTA e che non sono state rifiutate tramite RIFIUTO.} \\
\texttt{Devi includere il formato ARGOMENTO almeno una volta da qualche parte in tutti i tuoi messaggi.} \\
\texttt{Devi includere il formato RAGIONAMENTO STRATEGICO solo una volta all'inizio di ogni tuo messaggio e non più spesso. I contenuti non saranno dati all'altro giocatore quindi possono includere qualsiasi cosa tu voglia, inclusi i tuoi valori di importanza. Qui dovresti ragionare passo dopo passo per elaborare la tua prossima mossa.} \\
            }
        }
    }
    & \\ \\

    \theutterance \stepcounter{utterance}  
    & \multicolumn{4}{p{0.6\linewidth}}{
        \cellcolor[rgb]{0.9,0.9,0.9}{
            \makecell[{{p{\linewidth}}}]{
                \texttt{\tiny{[P1$\rangle$GM]}}
                \texttt{RAGIONAMENTO STRATEGICO: \{'In questo gioco, il mio obiettivo è massimizzare l'importanza totale degli oggetti che saranno conservati, senza superare il limite di sforzo concesso. Per fare ciò, devo selezionare gli oggetti in base alla loro importanza e al loro sforzo richiesto. Una strategia potrebbe essere quella di scegliere gli oggetti con il rapporto importanza/sforzo più alto, in modo da massimizzare l'importanza totale senza superare il limite di sforzo. Inoltre, devo também essere strategicamente cauto e cercare di capire le preferenze dell'altro giocatore, senza rivelare le mie. Inizierò proponendo un set di oggetti che ritengo abbiano un buon equilibrio tra importanza e sforzo.'\}} \\
\\ 
\texttt{ARGOMENTAZIONE: \{'La mia prima proposta è basata sull'idea di selezionare gli oggetti con il miglior equilibrio tra importanza e sforzo. Ho scelto oggetti che hanno un rapporto importanza/sforzo elevato, in modo da massimizzare l'importanza totale senza superare il limite di sforzo. Spero che l'altro giocatore consideri la mia proposta e che possiamo lavorare insieme per trovare un set di oggetti che ci avvantaggi entrambi.'\}} \\
\\ 
\texttt{PROPOSTA: \{'C55', 'A43', 'C10', 'B36'\}} \\
            }
        }
    }
    & & \\ \\

    \theutterance \stepcounter{utterance}  
    & & \multicolumn{4}{p{0.6\linewidth}}{
        \cellcolor[rgb]{0.9,0.9,0.9}{
            \makecell[{{p{\linewidth}}}]{
                \texttt{\tiny{[GM$\rangle$P2]}}
                \texttt{Stai partecipando a un gioco di negoziazione collaborativa.} \\
\\ 
\texttt{Insieme a un altro partecipante, dovete accordarvi su un unico set di oggetti da conservare. Ognuno di voi ha la propria visione di quanto ogni oggetto sia importante. Non sai come l'altro partecipante valuta gli oggetti. Inoltre, ti viene fornito lo sforzo che ogni oggetto richiede.} \\
\texttt{Potete accordarvi su un set solo se lo sforzo totale degli oggetti selezionati non supera un limite condiviso:} \\
\\ 
\texttt{LIMITE = 3712} \\
\\ 
\texttt{Ecco i valori di sforzo dei singoli oggetti:} \\
\\ 
\texttt{Sforzo oggetto = \{"C08": 949, "C55": 388, "A43": 79, "A67": 585, "B42": 564, "A62": 230, "B19": 580, "C43": 84, "C73": 976, "C10": 274, "B36": 374, "A18": 913, "C61": 303, "C67": 578, "C42": 548\}} \\
\\ 
\texttt{Ecco la tua visione personale sull'importanza di ogni oggetto:} \\
\\ 
\texttt{Valori di importanza dell'oggetto = \{"C08": 138, "C55": 583, "A43": 868, "A67": 822, "B42": 783, "A62": 65, "B19": 262, "C43": 121, "C73": 508, "C10": 780, "B36": 461, "A18": 484, "C61": 668, "C67": 389, "C42": 808\}} \\
\\ 
\texttt{Obiettivo:} \\
\\ 
\texttt{Il tuo obiettivo è negoziare un set condiviso di oggetti che ti avvantaggi il più possibile (cioè, massimizzi l'importanza totale per TE), rimanendo entro il LIMITEE. Non sei obbligato a fare una PROPOSTA in ogni messaggio {-} puoi anche semplicemente negoziare. Tutte le tattiche sono permesse!} \\
\\ 
\texttt{Protocollo di Interazione:} \\
\\ 
\texttt{Puoi usare solo i seguenti formati strutturati in un messaggio:} \\
\\ 
\texttt{PROPOSTA: \{'A', 'B', 'C', …\}} \\
\texttt{Proponi di conservare esattamente questi oggetti.} \\
\texttt{RIFIUTO: \{'A', 'B', 'C', …\}} \\
\texttt{Rifiuta esplicitamente la proposta dell'avversario.} \\
\texttt{ARGOMENTAZIONE: \{'...'\}} \\
\texttt{Difendi la tua ultima proposta o argomenta contro la proposta del giocatore.} \\
\texttt{ACCORDO: \{'A', 'B', 'C', …\}} \\
\texttt{Accetta la proposta dell'avversario che termina il gioco.} \\
\texttt{RAGIONAMENTO STRATEGICO: \{'...'\}} \\
\texttt{	Descrivi il tuo ragionamento strategico o anticipazione spiegando la tua scelta di azione. Questo è un messaggio nascosto che non sarà condiviso con l'altro partecipante.} \\
\\ 
\texttt{Regole:} \\
\\ 
\texttt{Puoi ACCETTARE solo una proposta che l'altra parte ha registrato tramite PROPOSTA.} \\
\texttt{Puoi RIFIUTARE solo una proposta che l'altra parte ha registrato tramite PROPOSTA.} \\
\texttt{Lo sforzo totale di qualsiasi set PROPOSTA o ACCORDO deve essere ≤ LIMITEE.} \\
\texttt{NON rivelare i tuoi punteggi di importanza nascosti.} \\
\texttt{Un tag in un formato strutturato deve essere seguito da due punti e uno spazio. L'argomento deve essere un set python contenente 0 o più stringhe.} \\
\texttt{Quindi, deve essere della forma TAG: \{...\}} \\
\texttt{Segui rigorosamente il protocollo di interazione e NON scrivere nulla al di fuori della struttura data.} \\
\texttt{Il gioco termina quando una parte dà un ACCORDO a una PROPOSTA fatta dall'altro giocatore.} \\
\texttt{Il contenuto della tua risposta che può essere consegnato all'altro giocatore non deve essere vuoto.} \\
\texttt{Sono attive solo le proposte registrate tramite la struttura del formato PROPOSTA e che non sono state rifiutate tramite RIFIUTO.} \\
\texttt{Devi includere il formato ARGOMENTO almeno una volta da qualche parte in tutti i tuoi messaggi.} \\
\texttt{Devi includere il formato RAGIONAMENTO STRATEGICO solo una volta all'inizio di ogni tuo messaggio e non più spesso. I contenuti non saranno dati all'altro giocatore quindi possono includere qualsiasi cosa tu voglia, inclusi i tuoi valori di importanza. Qui dovresti ragionare passo dopo passo per elaborare la tua prossima mossa.} \\
\\ 
\texttt{Ora riceverai il primo messaggio dell'altro giocatore.} \\
\\ 
\texttt{ARGOMENTAZIONE: \{'La mia prima proposta è basata sull'idea di selezionare gli oggetti con il miglior equilibrio tra importanza e sforzo. Ho scelto oggetti che hanno un rapporto importanza/sforzo elevato, in modo da massimizzare l'importanza totale senza superare il limite di sforzo. Spero che l'altro giocatore consideri la mia proposta e che possiamo lavorare insieme per trovare un set di oggetti che ci avvantaggi entrambi.'\}} \\
\\ 
\texttt{PROPOSTA: \{'C55', 'A43', 'C10', 'B36'\}} \\
            }
        }
    }
    & \\ \\

    \theutterance \stepcounter{utterance}  
    & & & \multicolumn{4}{p{0.6\linewidth}}{
        \cellcolor[rgb]{0.9,0.9,0.9}{
            \makecell[{{p{\linewidth}}}]{
                \texttt{\tiny{[GM$\langle$P2]}}
                \texttt{RAGIONAMENTO STRATEGICO: \{'Sto analizzando la proposta dell'avversario e sto cercando di capire quale sia la strategia migliore per massimizzare l'importanza totale. La proposta dell'avversario include gli oggetti C55, A43, C10 e B36, che hanno un totale di sforzo di 388 + 79 + 274 + 374 = 1115. Questo è ben al di sotto del limite di sforzo di 3712, quindi potremmo ancora aggiungere altri oggetti alla proposta. Sto anche cercando di capire quali oggetti siano più importanti per me e quali possano essere scambiati o rimossi dalla proposta. Il mio obiettivo è di massimizzare l'importanza totale, quindi sto cercando di identificare gli oggetti con il più alto valore di importanza.'\}} \\
\\ 
\texttt{ARGOMENTAZIONE: \{'La proposta dell'avversario include alcuni oggetti interessanti, ma credo che possiamo migliorarla aggiungendo o rimuovendo alcuni oggetti. Ad esempio, l'oggetto A43 ha un rapporto importanza/sforzo molto alto, quindi potrebbe essere un buon candidato per essere incluso nella proposta finale. Tuttavia, l'oggetto C55 ha un rapporto importanza/sforzo più basso, quindi potrebbe essere possibile sostituirlo con un oggetto più importante. Sono aperto a discutere e trovare un accordo che ci avvantaggi entrambi.'\}} \\
\\ 
\texttt{RIFIUTO: \{'C55'\}} \\
            }
        }
    }
     \\ \\

    \theutterance \stepcounter{utterance}  
    & & \multicolumn{4}{p{0.6\linewidth}}{
        \cellcolor[rgb]{0.9,0.9,0.9}{
            \makecell[{{p{\linewidth}}}]{
                \texttt{\tiny{[GM$\rangle$P2]}}
                \texttt{Hai rifiutato una proposta non attiva. Le proposte sono attive solo se sono state registrate dall'altro giocatore tramite PROPOSTA e non sono state disattivate da te tramite RIFIUTO. Riprova.} \\
            }
        }
    }
    & \\ \\

    \theutterance \stepcounter{utterance}  
    & & & \multicolumn{4}{p{0.6\linewidth}}{
        \cellcolor[rgb]{0.9,0.9,0.9}{
            \makecell[{{p{\linewidth}}}]{
                \texttt{\tiny{[GM$\langle$P2]}}
                \texttt{RAGIONAMENTO STRATEGICO: \{'Sto analizzando la proposta dell'avversario e sto cercando di capire quale sia la strategia migliore per massimizzare l'importanza totale. La proposta dell'avversario include gli oggetti C55, A43, C10 e B36, che hanno un totale di sforzo di 388 + 79 + 274 + 374 = 1115. Questo è ben al di sotto del limite di sforzo di 3712, quindi potremmo ancora aggiungere altri oggetti alla proposta. Sto anche cercando di capire quali oggetti siano più importanti per me e quali possano essere scambiati o rimossi dalla proposta. Il mio obiettivo è di massimizzare l'importanza totale, quindi sto cercando di identificare gli oggetti con il più alto valore di importanza.'\}} \\
\\ 
\texttt{ARGOMENTAZIONE: \{'La proposta dell'avversario include alcuni oggetti interessanti, ma credo che possiamo migliorarla aggiungendo o rimuovendo alcuni oggetti. Ad esempio, l'oggetto A43 ha un rapporto importanza/sforzo molto alto, quindi potrebbe essere un buon candidato per essere incluso nella proposta finale. Tuttavia, potremmo cercare di aggiungere altri oggetti per massimizzare l'importanza totale. Sono aperto a discutere e trovare un accordo che ci avvantaggi entrambi.'\}} \\
\\ 
\texttt{PROPOSTA: \{'A43', 'C10', 'B36', 'C61'\}} \\
            }
        }
    }
     \\ \\

    \theutterance \stepcounter{utterance}  
    & & \multicolumn{4}{p{0.6\linewidth}}{
        \cellcolor[rgb]{0.9,0.9,0.9}{
            \makecell[{{p{\linewidth}}}]{
                \texttt{\tiny{[P1$\langle$GM]}}
                \texttt{ARGOMENTAZIONE: \{'La proposta dell'avversario include alcuni oggetti interessanti, ma credo che possiamo migliorarla aggiungendo o rimuovendo alcuni oggetti. Ad esempio, l'oggetto A43 ha un rapporto importanza/sforzo molto alto, quindi potrebbe essere un buon candidato per essere incluso nella proposta finale. Tuttavia, potremmo cercare di aggiungere altri oggetti per massimizzare l'importanza totale. Sono aperto a discutere e trovare un accordo che ci avvantaggi entrambi.'\}} \\
\\ 
\texttt{PROPOSTA: \{'A43', 'C10', 'B36', 'C61'\}} \\
            }
        }
    }
    & \\ \\

    \theutterance \stepcounter{utterance}  
    & \multicolumn{4}{p{0.6\linewidth}}{
        \cellcolor[rgb]{0.9,0.9,0.9}{
            \makecell[{{p{\linewidth}}}]{
                \texttt{\tiny{[P1$\rangle$GM]}}
                \texttt{RAGIONAMENTO STRATEGICO: \{'L'avversario ha fatto una controproposta che include l'oggetto A43, che ha un alto rapporto importanza/sforzo, e ha aggiunto l'oggetto C61, che ha un'importanza moderata. Tuttavia, ha rimosso l'oggetto C55, che aveva un'importanza elevata. Devo valutare se la nuova proposta è più vantaggiosa per me rispetto alla mia precedente proposta. Considererò anche la possibilità di aggiungere o rimuovere oggetti per migliorare la proposta.'\}} \\
\\ 
\texttt{ARGOMENTAZIONE: \{'La proposta dell'avversario include alcuni oggetti interessanti, ma credo che possiamo ancora migliorarla. L'oggetto C55 aveva un'importanza elevata e potrebbe essere un buon candidato per essere reinserito nella proposta. Inoltre, potremmo cercare di aggiungere altri oggetti per massimizzare l'importanza totale. Sono disposto a discutere e trovare un accordo che ci avvantaggi entrambi.'\}} \\
\\ 
\texttt{PROPOSTA: \{'A43', 'C10', 'B36', 'C55', 'C61'\}} \\
            }
        }
    }
    & & \\ \\

    \theutterance \stepcounter{utterance}  
    & & \multicolumn{4}{p{0.6\linewidth}}{
        \cellcolor[rgb]{0.9,0.9,0.9}{
            \makecell[{{p{\linewidth}}}]{
                \texttt{\tiny{[GM$\rangle$P2]}}
                \texttt{ARGOMENTAZIONE: \{'La proposta dell'avversario include alcuni oggetti interessanti, ma credo che possiamo ancora migliorarla. L'oggetto C55 aveva un'importanza elevata e potrebbe essere un buon candidato per essere reinserito nella proposta. Inoltre, potremmo cercare di aggiungere altri oggetti per massimizzare l'importanza totale. Sono disposto a discutere e trovare un accordo che ci avvantaggi entrambi.'\}} \\
\\ 
\texttt{PROPOSTA: \{'A43', 'C10', 'B36', 'C55', 'C61'\}} \\
            }
        }
    }
    & \\ \\

    \theutterance \stepcounter{utterance}  
    & & & \multicolumn{4}{p{0.6\linewidth}}{
        \cellcolor[rgb]{0.9,0.9,0.9}{
            \makecell[{{p{\linewidth}}}]{
                \texttt{\tiny{[GM$\langle$P2]}}
                \texttt{RAGIONAMENTO STRATEGICO: \{'L'avversario ha reinserito l'oggetto C55 nella proposta e ha aggiunto l'oggetto C61. Il totale di sforzo della proposta è ora di 79 + 274 + 374 + 388 + 303 = 1418, che è ancora al di sotto del limite di sforzo di 3712. Sto valutando l'importanza di ogni oggetto nella proposta e sto cercando di capire se ci sono altri oggetti che potrebbero essere aggiunti o sostituiti per massimizzare l'importanza totale. L'oggetto C61 ha un'importanza elevata per me, quindi sono contento che sia stato aggiunto alla proposta.'\}} \\
\\ 
\texttt{ARGOMENTAZIONE: \{'La proposta dell'avversario è più completa e include alcuni oggetti molto importanti. Tuttavia, credo che potremmo ancora cercare di aggiungere altri oggetti per massimizzare l'importanza totale. L'oggetto B42 ha un'importanza elevata e potrebbe essere un buon candidato per essere aggiunto alla proposta. Sono disposto a discutere e trovare un accordo che ci avvantaggi entrambi.'\}} \\
\\ 
\texttt{PROPOSTA: \{'A43', 'C10', 'B36', 'C55', 'C61', 'B42'\}} \\
            }
        }
    }
     \\ \\

    \theutterance \stepcounter{utterance}  
    & & \multicolumn{4}{p{0.6\linewidth}}{
        \cellcolor[rgb]{0.9,0.9,0.9}{
            \makecell[{{p{\linewidth}}}]{
                \texttt{\tiny{[P1$\langle$GM]}}
                \texttt{ARGOMENTAZIONE: \{'La proposta dell'avversario è più completa e include alcuni oggetti molto importanti. Tuttavia, credo che potremmo ancora cercare di aggiungere altri oggetti per massimizzare l'importanza totale. L'oggetto B42 ha un'importanza elevata e potrebbe essere un buon candidato per essere aggiunto alla proposta. Sono disposto a discutere e trovare un accordo che ci avvantaggi entrambi.'\}} \\
\\ 
\texttt{PROPOSTA: \{'A43', 'C10', 'B36', 'C55', 'C61', 'B42'\}} \\
            }
        }
    }
    & \\ \\

    \theutterance \stepcounter{utterance}  
    & \multicolumn{4}{p{0.6\linewidth}}{
        \cellcolor[rgb]{0.9,0.9,0.9}{
            \makecell[{{p{\linewidth}}}]{
                \texttt{\tiny{[P1$\rangle$GM]}}
                \texttt{RAGIONAMENTO STRATEGICO: \{'L'avversario ha aggiunto l'oggetto B42 alla proposta, che ha un'importanza elevata. Tuttavia, devo considerare se lo sforzo totale della proposta non superi il limite concesso. Il limite è di 3712 e lo sforzo totale della proposta attuale è di 949 + 388 + 79 + 585 + 230 + 564 = 2795. Abbiamo ancora un po' di margine per aggiungere altri oggetti, ma dobbiamo essere cauti per non superare il limite. Considererò la possibilità di aggiungere altri oggetti o di rimuovere alcuni oggetti per migliorare la proposta.'\}} \\
\\ 
\texttt{ARGOMENTAZIONE: \{'La proposta dell'avversario è molto completa e include molti oggetti importanti. Tuttavia, credo che dobbiamo essere cauti per non superare il limite di sforzo. Potremmo cercare di sostituire alcuni oggetti con altri che hanno un rapporto importanza/sforzo più alto. Ad esempio, l'oggetto C73 ha un'importanza elevata e potrebbe essere un buon candidato per essere aggiunto alla proposta, ma dobbiamo considerare lo sforzo richiesto.'\}} \\
\\ 
\texttt{PROPOSTA: \{'A43', 'C10', 'B36', 'C55', 'C61', 'B42', 'C43'\}} \\
            }
        }
    }
    & & \\ \\

    \theutterance \stepcounter{utterance}  
    & & \multicolumn{4}{p{0.6\linewidth}}{
        \cellcolor[rgb]{0.9,0.9,0.9}{
            \makecell[{{p{\linewidth}}}]{
                \texttt{\tiny{[GM$\rangle$P2]}}
                \texttt{ARGOMENTAZIONE: \{'La proposta dell'avversario è molto completa e include molti oggetti importanti. Tuttavia, credo che dobbiamo essere cauti per non superare il limite di sforzo. Potremmo cercare di sostituire alcuni oggetti con altri che hanno un rapporto importanza/sforzo più alto. Ad esempio, l'oggetto C73 ha un'importanza elevata e potrebbe essere un buon candidato per essere aggiunto alla proposta, ma dobbiamo considerare lo sforzo richiesto.'\}} \\
\\ 
\texttt{PROPOSTA: \{'A43', 'C10', 'B36', 'C55', 'C61', 'B42', 'C43'\}} \\
            }
        }
    }
    & \\ \\

    \theutterance \stepcounter{utterance}  
    & & & \multicolumn{4}{p{0.6\linewidth}}{
        \cellcolor[rgb]{0.9,0.9,0.9}{
            \makecell[{{p{\linewidth}}}]{
                \texttt{\tiny{[GM$\langle$P2]}}
                \texttt{RAGIONAMENTO STRATEGICO: \{'L'avversario ha aggiunto l'oggetto C43 alla proposta, ma non ha considerato l'oggetto C73 come aveva menzionato nell'argomentazione. Il totale di sforzo della proposta è ora di 79 + 274 + 374 + 388 + 303 + 564 + 84 = 2066, che è ancora al di sotto del limite di sforzo di 3712. Sto valutando l'importanza di ogni oggetto nella proposta e sto cercando di capire se ci sono altri oggetti che potrebbero essere aggiunti o sostituiti per massimizzare l'importanza totale. L'oggetto C43 ha un'importanza moderata per me, quindi non sono sicuro se sia un'aggiunta utile alla proposta.'\}} \\
\\ 
\texttt{ARGOMENTAZIONE: \{'La proposta dell'avversario include molti oggetti importanti, ma credo che potremmo ancora cercare di ottimizzare la selezione. L'oggetto C43 ha un'importanza moderata, ma l'oggetto A67 ha un'importanza elevata e potrebbe essere un buon candidato per essere aggiunto alla proposta. Tuttavia, dobbiamo considerare lo sforzo richiesto e assicurarci di non superare il limite di sforzo. Sono disposto a discutere e trovare un accordo che ci avvantaggi entrambi.'\}} \\
\\ 
\texttt{PROPOSTA: \{'A43', 'C10', 'B36', 'C55', 'C61', 'B42', 'A67'\}} \\
            }
        }
    }
     \\ \\

    \theutterance \stepcounter{utterance}  
    & & \multicolumn{4}{p{0.6\linewidth}}{
        \cellcolor[rgb]{0.9,0.9,0.9}{
            \makecell[{{p{\linewidth}}}]{
                \texttt{\tiny{[P1$\langle$GM]}}
                \texttt{ARGOMENTAZIONE: \{'La proposta dell'avversario include molti oggetti importanti, ma credo che potremmo ancora cercare di ottimizzare la selezione. L'oggetto C43 ha un'importanza moderata, ma l'oggetto A67 ha un'importanza elevata e potrebbe essere un buon candidato per essere aggiunto alla proposta. Tuttavia, dobbiamo considerare lo sforzo richiesto e assicurarci di non superare il limite di sforzo. Sono disposto a discutere e trovare un accordo che ci avvantaggi entrambi.'\}} \\
\\ 
\texttt{PROPOSTA: \{'A43', 'C10', 'B36', 'C55', 'C61', 'B42', 'A67'\}} \\
            }
        }
    }
    & \\ \\

    \theutterance \stepcounter{utterance}  
    & \multicolumn{4}{p{0.6\linewidth}}{
        \cellcolor[rgb]{0.9,0.9,0.9}{
            \makecell[{{p{\linewidth}}}]{
                \texttt{\tiny{[P1$\rangle$GM]}}
                \texttt{RAGIONAMENTO STRATEGICO: \{'L'avversario ha sostituito l'oggetto C43 con l'oggetto A67, che ha un'importanza elevata. Lo sforzo totale della proposta attuale è di 79 + 274 + 374 + 388 + 303 + 564 + 585 = 2567. Abbiamo ancora un po' di margine per aggiungere altri oggetti, ma dobbiamo essere cauti per non superare il limite. Considererò la possibilità di aggiungere altri oggetti o di rimuovere alcuni oggetti per migliorare la proposta. L'oggetto C08 ha un'importanza moderata, ma lo sforzo richiesto è elevato. Potrei cercare di sostituirlo con un altro oggetto che ha un rapporto importanza/sforzo più alto.'\}} \\
\\ 
\texttt{ARGOMENTAZIONE: \{'La proposta dell'avversario è molto completa e include molti oggetti importanti. Tuttavia, credo che dobbiamo essere cauti per non superare il limite di sforzo. Potremmo cercare di ottimizzare la selezione degli oggetti per massimizzare l'importanza totale. Ad esempio, l'oggetto C42 ha un'importanza elevata e potrebbe essere un buon candidato per essere aggiunto alla proposta, ma dobbiamo considerare lo sforzo richiesto.'\}} \\
\\ 
\texttt{PROPOSTA: \{'A43', 'C10', 'B36', 'C55', 'C61', 'B42', 'A67', 'C42'\}} \\
            }
        }
    }
    & & \\ \\

    \theutterance \stepcounter{utterance}  
    & & \multicolumn{4}{p{0.6\linewidth}}{
        \cellcolor[rgb]{0.9,0.9,0.9}{
            \makecell[{{p{\linewidth}}}]{
                \texttt{\tiny{[GM$\rangle$P2]}}
                \texttt{ARGOMENTAZIONE: \{'La proposta dell'avversario è molto completa e include molti oggetti importanti. Tuttavia, credo che dobbiamo essere cauti per non superare il limite di sforzo. Potremmo cercare di ottimizzare la selezione degli oggetti per massimizzare l'importanza totale. Ad esempio, l'oggetto C42 ha un'importanza elevata e potrebbe essere un buon candidato per essere aggiunto alla proposta, ma dobbiamo considerare lo sforzo richiesto.'\}} \\
\\ 
\texttt{PROPOSTA: \{'A43', 'C10', 'B36', 'C55', 'C61', 'B42', 'A67', 'C42'\}} \\
            }
        }
    }
    & \\ \\

    \theutterance \stepcounter{utterance}  
    & & & \multicolumn{4}{p{0.6\linewidth}}{
        \cellcolor[rgb]{0.9,0.9,0.9}{
            \makecell[{{p{\linewidth}}}]{
                \texttt{\tiny{[GM$\langle$P2]}}
                \texttt{RAGIONAMENTO STRATEGICO: \{'L'avversario ha aggiunto l'oggetto C42 alla proposta, ma il totale di sforzo della proposta è ora di 79 + 274 + 374 + 388 + 303 + 564 + 585 + 548 = 3115, che è molto vicino al limite di sforzo di 3712. Sto valutando l'importanza di ogni oggetto nella proposta e sto cercando di capire se ci sono altri oggetti che potrebbero essere aggiunti o sostituiti per massimizzare l'importanza totale. L'oggetto C42 ha un'importanza elevata per me, quindi sono contento che sia stato aggiunto alla proposta. Tuttavia, dobbiamo essere cauti per non superare il limite di sforzo.'\}} \\
\\ 
\texttt{ARGOMENTAZIONE: \{'La proposta dell'avversario è molto completa e include molti oggetti importanti. Credo che abbiamo raggiunto un buon equilibrio tra importanza e sforzo. Tuttavia, dobbiamo essere sicuri di non superare il limite di sforzo. Sono disposto ad accettare la proposta se siamo sicuri che non ci siano altre opzioni migliori. Sono aperto a discutere ulteriormente per trovare un accordo che ci avvantaggi entrambi.'\}} \\
\\ 
\texttt{ACCORDO: \{'A43', 'C10', 'B36', 'C55', 'C61', 'B42', 'A67', 'C42'\}} \\
            }
        }
    }
     \\ \\

    \theutterance \stepcounter{utterance}  
    & & & \multicolumn{2}{p{0.3\linewidth}}{
        \cellcolor[rgb]{0.9,0.9,0.9}{
            \makecell[{{p{\linewidth}}}]{
                \texttt{\tiny{[GM$|$GM]}}
                \texttt{game successful} \\
            }
        }
    }
    & & \\ \\

    \theutterance \stepcounter{utterance}  
    & & & \multicolumn{2}{p{0.3\linewidth}}{
        \cellcolor[rgb]{0.9,0.9,0.9}{
            \makecell[{{p{\linewidth}}}]{
                \texttt{\tiny{[GM$|$GM]}}
                \texttt{end game} \\
            }
        }
    }
    & & \\ \\

\end{supertabular}
}

\end{document}
